\section{Experimental setup}
\label{sec:setup}

\subsection{Training data}
\label{sec:data}

We follow LLM-as-Jev \citep{li2026llmasjev} and train on its Reason+Long mixture, with Nimble added.
It has 13 datasets: intent routing (CLINC150, MASSIVE, Bitext), short reasoning (ReClor, LogiQA~2.0, CosmosQA, PIQA, ANLI, ProofWriter), long input (ContractNLI, MuSiQue, QuALITY) and Nimble.
Each is turned into a state (the passages), a question and a list of options.
We keep \NPerDataset{} training rows per dataset, \NTrain{} in total.
Rows are filtered in this order:
\begin{enumerate}\setlength{\itemsep}{0pt}\setlength{\topsep}{2pt}
  \item \textbf{Duplicates.} Rows with the same input are collapsed to one, and rows with the same input but different answers are removed.
  \item \textbf{Leakage between splits.} A training row is removed if its input, or the source document it comes from, also appears in the validation split.
  \item \textbf{Benchmark overlap.} A training row is removed if it shares a 13-word n-gram, or an identical sentence of six or more words, with any evaluation benchmark.
  \item \textbf{Length.} Prompts longer than \MaxTokens{} tokens are dropped, never truncated.
\end{enumerate}
The held-out rows are drawn from the validation split with the same duplicate and length filters, up to 500 per dataset and in whole source documents, giving \NHeldOut{} rows.
We split them into a dev set of \NDev{} rows, used to choose learning rates, and a test set of \NTest{} rows that no choice depends on.

\subsection{Models}
\label{sec:models}

\paragraph{Encoder.}
UltraBERT-380M \citep{ultrabert} is fully fine-tuned with the format of \cref{sec:format}, once with the plain softmax and once with the sigmoid-normalised loss.

\paragraph{Decoders.}
We train Qwen3-0.6B \citep{qwen3} with the full LLM-as-Jev recipe (trie loss and the three anchors, $\lambda=1$), updating all weights.
The anchors use a frozen copy of the starting model as $p_0$.
We also report two LoRA models \citep{lora} that we trained earlier with the same recipe and released, Qwen3-0.6B and Qwen3-4B.
They reach \TestQwLora{}\% and \TestQwFourLora{}\% on the test set (\cref{tab:headline}).
Full fine-tuning of the 0.6B decoder reaches \TestQwFull{}\%, \DiffQwFullQwLora{} points above LoRA, and the encoder is fully fine-tuned, so full fine-tuning is the matched comparison and the LoRA models are references.
They were trained on the 12K-row release of the data without ReClor, while all new runs use the 13K-row mixture with ReClor.

\subsection{Protocol}
\label{sec:protocol}

Every model is trained for one epoch with a batch of 32 rows.
We sweep the learning rate with one seed per value and choose by accuracy on the dev set:
$\{8\cdot10^{-6}, 1.5\cdot10^{-5}, 3\cdot10^{-5}, 5\cdot10^{-5}\}$ for the encoder with the softmax, $\{8\cdot10^{-6}, 1.5\cdot10^{-5}, 2\cdot10^{-5}, 3\cdot10^{-5}\}$ with the sigmoid-normalised loss, and $\{5\cdot10^{-6}, 10^{-5}, 2\cdot10^{-5}\}$ for the decoder.
The two best Qwen rates differ by 0.1 points on dev, and we use $10^{-5}$.
At the chosen rate we repeat the run with different seeds, which change the order of the rows and the option shuffling, and report the average.
The test set is not used for any of these choices.

The metric is accuracy.
ALL is the mean accuracy over the 13 datasets, so that every dataset counts equally.
On the test set we average each row's correctness over the seeds of a model and report ALL of this average.
The encoder is scored with fp32 weights under bf16 autocast, because bf16 weights change its long-document results.
The Qwen checkpoints are scored with bf16 weights, as the released models are.

\subsection{External benchmarks and cost}
\label{sec:external}

To see whether training on the mixture changes behaviour outside it, we score six benchmarks as decision rows: BoolQ, MMLU, MMLU-Pro, ARC-Challenge, Winogrande and SciQ, with 500, 1200, 800, 600, 500 and 400 rows.
The sample sizes follow LLM-as-Jev, whose samples are not released, so ours are seeded uniform draws.
External is the mean accuracy over the six.
These benchmarks are not used for any choice.

Time per decision is the mean over 1{,}300 held-out rows on one A100, with one row at a time.
